\documentclass[12pt,a4paper]{scrartcl}
\usepackage[utf8]{inputenc}
\usepackage[T1]{fontenc}
\usepackage{ngerman}
\usepackage{placeins}
%\usepackage[ngerman]{babel}  % use babel for multi-ling or english
%
%\usepackage{amsmath,amssymb,amstext}

\usepackage{amsmath}
\usepackage{multirow}
%\usepackage{sv}
\usepackage{pgf}
\usepackage{pst-sigsys}
\usepackage{multido}
\usepackage{pdfpages}
\usepackage{listings}
\usepackage{color}
\definecolor{light-gray}{gray}{0.95}


% fuer Zitate
\usepackage[numbers]{natbib}
\usepackage{nomencl,longtable} 
% Festlegung Art der Zitierung - Havardmethode: Abkuerzung Autor + Jahr
%\bibliographystyle{alphadin}

\bibliographystyle{unsrt}
%\usepackage[twoside=false,
%            left=2cm,
%            right=2cm,
%            top=2cm,
%            bottom=3cm]{geometry}
\usepackage{tikz}
\usepackage{siunitx}

\usetikzlibrary{shapes,arrows,automata,backgrounds,calendar}
\usepackage[europeancurrents,
						europeanvoltages,
						europeanresistors,
						cuteinductors,
						europeanports]{circuitikz}
\usepgflibrary{plotmarks}
\usepackage{trfsigns}
\usepackage{pstricks} \SpecialCoor
\makeatletter
\def\psStartPunkt(#1){\pst@getcoor{#1}\pst@tempA 
  \pstVerb{\pst@tempA 
           \pst@number\psyunit div /cp.Y exch def 
           \pst@number\psxunit div /cp.X exch def }}
\def\psVektor{\pst@object{psVektor}}
\def\psVektor@i(#1){%
  \pst@killglue%
  \pst@getcoor{#1}\pst@tempA%
  \begin@SpecialObj%
  \rput(! cp.X cp.Y ){\psline[arrowsize=6pt]{->}(0,0)(#1)%
    \psarc[arrows=->,arrowsize=4pt](0,0){1}{0}{!\pst@tempA exch atan}%
    \psline[linestyle=dotted](1.5,0)}%
  \end@SpecialObj%
  \pstVerb{tx@Dict begin \pst@tempA \pst@number\psyunit div 
            cp.Y add /cp.Y exch def 
            \pst@number\psxunit div 
            cp.X add /cp.X exch def end}  
  \ignorespaces%
}
%%%%%%%%%%%%%%%%%%%%%%%%%%%%%%%%%%%%%%%%%%%%%%%%%%%%%%%%%%%%%%%%%%%%%%%%%%%%%%%%
%%% Headings
%%%
%\usepackage[automark]{scrpage2}
%\pagestyle{scrheadings}
%\ihead[]{}
%\chead[]{}
%\ohead[]{}
%\ifoot[]{}
%\cfoot[]{}
%\ofoot[]{}

%%%%%%%%%%%%%%%%%%%%%%%%%%%%%%%%%%%%%%%%%%%%%%%%%%%%%%%%%%%%%%%%%%%%%%%%%%%%%%%%
%%% Commands
%%%
\renewcommand{\vec}[1]{\boldsymbol{#1}}
\newcommand{\mat}[1]{\boldsymbol{#1}}

\DeclareMathAlphabet{\mathpzc}{OT1}{pzc}{m}{it}
\newcommand{\z}{\mathpzc{z}}

\newcommand{\nomunit}[1]{% 
\renewcommand{\nomentryend}{\hspace{2em}\hspace*{\fill}#1}}
%%%%%%%%%%%%%%%%%%%%%%%%%%%%%%%%%%%%%%%%%%%%%%%%%%%%%%%%%%%%%%%%%%%%%%


%\title{Übung Drehfeldmaschinen \\mit Umrichter}

\author{}

%\date{Graz, am \today{}}

\parindent0pt

%%%%%%%%%%%%%%%%%%%%%%%%%%%%%% Main %%%%%%%%%%%%%%%%%%%%%%%%%%%%%%%%%% 


\begin{document}
\begin{titlepage}
\pagestyle{empty} \enlargethispage*{25cm}\samepage{

\vspace*{-1.5cm}
\begin{center}
\begin{minipage}[!h]{18cm}
\hspace*{-1.1cm}
\includegraphics[width=3.3cm]{igte}
\begin{tabular}{p{10cm}}
\centering{
\Large Institut für Grundlagen und Theorie\\
der Elektrotechnik\\
Technische Universität Graz
}
\end{tabular}
\includegraphics[width=3.3cm]{TUG}
\end{minipage}

%\vspace*{2.4cm}
\vspace*{5cm}
%
\Huge {Seminararbeit:\\
Finite Elemente Software \\zur Lösung von \\elektrostatischen und stationären Strömungsfeld-Probleme} %
%\vspace*{.4cm} \Large{ Sommersemester 2015\\}
%%
%\vspace*{2.5cm} 
\vspace*{2cm}
%\Huge{\textbf{Übung 4 – Kurzschluss 2}\\}\vspace*{.7cm} 
%%
%\Large{Gruppe: 10\\} \vspace*{.7cm}%
%\Large{Teilnehmer:\\} \vspace*{.1cm}
% Name der beteiligten Studierenden

\Large{vorgelegt von:\\
Tobias Florian Lafer (1530012)\\
am \today} 
\vspace{9cm}


%
%
\Large{Betreuer: Dipl.-Ing. Dr.techn. Thomas Bauernfeind}\vspace*{2cm}\vfill
%
%\Large{Übungsdatum: 02.06.2015}
\end{center}}%
\clearpage 

\tableofcontents
\newpage

\end{titlepage}
%



\newcommand{\screenshot}[3]{
\begin{figure}[ht]
 \centering
\hspace*{-.05\textwidth}
\makebox[\textwidth]{
  \includegraphics[width = 19cm]{#3}}
 \caption{#1} \label{fig:#2} 
\end{figure}
\FloatBarrier
}


\newcommand{\screenshotWidth}[4]{
\begin{figure}[ht]
 \centering
\hspace*{-.05\textwidth}
\makebox[\textwidth]{
  \includegraphics[width = {#4 cm}]{#3}}
 \caption{#1.} \label{fig:#2} 
\end{figure}
\FloatBarrier
}

\newcommand{\screenshotHeight}[4]{
\begin{figure}[ht]
 \centering
\hspace*{-.05\textwidth}
\makebox[\textwidth]{
  \includegraphics[height = {#4 cm}]{#3}}
 \caption{#1} \label{fig:#2} 
\end{figure}
\FloatBarrier
}

\newcommand{\doubleScreenShot}[9]{

  \begin{figure}[!ht]
	  \caption*{#9}
    \begin{minipage}{0.49\linewidth}\centering
      %\rule{\linewidth}{0.5\linewidth}  % Platzhalter
			\includegraphics[width={#4}\textwidth]{#3} 
      \caption{#1.}\label{fig:#2}
    \end{minipage}
    \hfill
    \begin{minipage}{0.49\linewidth}\centering
      %\rule{\linewidth}{0.5\linewidth}  % Platzhalter
			\includegraphics[width={#8}\textwidth]{#7}
      \caption{#5}\label{fig:#6}
    \end{minipage}
		
  \end{figure} 
\FloatBarrier
}




%Aufgabenstellung

%Analytisches Modell
%	gleichmäßige Stromverteilung
%	Sinusförmige Stromverteilung
%	reale Stromverteilung
%	Strom (Storer)
%	Werner
%	Diskussion
%FEM-Modell
%	Formulierung
%	Modell
%	dünner Feedgap
%	dicker Feedgap
%	Stromfaden
%Diskussion


\section*{Zusammenfassung}
\vspace*{2cm}
\begin{large}

Die Methode der finiten Elemente ist eine der bekanntesten und am weitesten Verbreiteten Methoden zur Lösung von Randwertproblemen in unterschiedlichsten Disziplinen. Das Institut für Grundlagen und Theorie der Elektrotechnik hat zu diesem Zwecke vor vielen Jahren mit der Entwicklung einer kommerziell Verwendbaren Software für eben jene Probleme auf dem Bereich der Elektrotechnik entwickelt. Aus diesen Bemühungen sind die Softwarepakete EleFAnT3D und EleFAnt2D entstanden.\newline

Beim Einsatz in der Lehre geht der Fokus einer solchen Software jedoch weg von hoher Optimierung und großem Funktionsumfang hin zu einfacher Bedienung und leichter Erweiterbarkeit des Source-Codes. Aus diesen Überlegungen heraus wurde die Entwicklung einer, auf der Open-Source-Software Gmsh und dem weit verbreiteten Mathematik-Programm Matlab\textregistered, basierenden Software zur Lösung oben genannter Probleme zum Einsatz speziell in der Lehre gestartet.\newline

Diese Seminararbeit versteht sich als ein Teil vergangener, laufender und noch folgender Seminararbeiten zur sukzessiven Entwicklung dieser Software.  \newline
Ziel dieser Seminararbeit ist die Implementierung der Lösung von zweidimensionalen, ebenen elektrostatischen und stationären Strömungsfeld-Problemen, sowie die automatisierte Generierung des Gitters mittels der Software Gmsh, dessen Import und Verarbeitung in der Software.

\end{large}


\section{Formfunktionen}
\subsection{Erster Ordnung}
\subsection{Zweiter Ordnung}
\subsection{Dritter Ordnung}
\section{Finite-Elemente}
\subsection{3-Punkte linear}
\subsection{6-Punkte quadratisch}
\subsection{10-Punkte kubisch}

\section{Gmsh}
\subsection{Erstellen der Geometrie}
\subsection{Generieren und exportieren des Gitters}

\section{Import der Geometrie und des Gitters}
\section{Assemblierung und Lösung des Gleichungssystems}

\section{Grafischer Postprozessor}



\newpage
\addcontentsline{toc}{section}{Literatur}
\bibliography{literatur}

\end{document}